% I. Introduction
Low back pain is the leading cause of years lived with disability worldwide: in 2020 it affected about 619 million people, and it is projected to reach 843 million cases by 2050 \cite{gbd2023}. Magnetic resonance imaging (MRI) is the reference imaging test for the intervertebral disc. On sagittal T2-weighted slices, loss of disc signal and height indicates degeneration, according to the Pfirrmann scale \cite{pfirrmann2001}, and the extension of disc material beyond its margins defines bulges and herniations, according to the consensus nomenclature \cite{fardon2014}.

Reading these studies is manual, time-consuming and dependent on the reader's experience. Even with a standardized nomenclature, inter-radiologist agreement on disc contour is only fair to moderate \cite{arana2011}. A support system that automatically classifies the state of each disc could shorten reading time and make criteria more consistent. Computer vision has already been applied to this task with good results \cite{sari2025,tan2024,sayed2025scirep,wang2025}, but many figures come from small test sets, from binary tasks or from splits not declared at the patient level, and they drop when evaluated on external data \cite{yi2023,nigru2024}. When the same patient contributes images to both training and test sets, the model may recognize the patient instead of the pathology, inflating the metrics \cite{brunekreef2024,rumala2023}.

This work presents DeepSpine, an end-to-end pipeline to classify lumbar discs (L1-L2 to L5-S1) on sagittal T2-weighted slices of the \emph{Lumbar Spine MRI Dataset} \cite{sudirman2019}. The initial project proposal defined four classes (Normal, Degeneration, Bulge and Herniation) and three research questions:
\begin{enumerate}
  \item[Q1.] \textbf{Exploratory analysis:} how are the cases of each class distributed in the dataset, and which visual differences (brightness, shape or size) distinguish a healthy disc from a pathological one?
  \item[Q2.] \textbf{Classical versus deep learning:} can simpler and faster classical algorithms (SVM or Random Forest) achieve a reliable diagnosis from image features, or is a deep neural network required?
  \item[Q3.] \textbf{Deep learning:} which convolutional neural network architecture achieves the highest accuracy (with a target above 90\%) and the lowest false-negative rate?
\end{enumerate}
During dataset construction, Herniation and Bulge proved inseparable on the sagittal plane and were merged into a single class, Protrusion (Section~\ref{sec:annotation}); the three questions are therefore answered over three classes.

The contributions are threefold: (i) a dataset of 5,123 ROIs from 399 patients, built with a traceable process that is reproducible from the public data and that declares its errors and limitations; (ii) a comparison of four models from different families (YOLOv8s-cls, ResNet-50, SVM and Random Forest) under the same patient-level split, the same metrics and evaluation on original images; and (iii) an analysis of class imbalance and vertebral-level bias that puts the achieved accuracy into perspective. Section~\ref{sec:related} reviews prior work; Section~\ref{sec:methodology} describes the dataset construction, the problem formulation and the models; Section~\ref{sec:results} presents and discusses the results; and Section~\ref{sec:conclusions} answers the research questions and outlines limitations and future work.
